% GENERATED FILE. Edit canonical lessons, then run npm run generate:books.
% canonical-source-hash: fnv1a64:30f7a17628e81e7a
% canonical-lessons: AR-C467-mahmiyya, AR-C467-muhaddad, AR-C467-hafaza, AR-C467-inqarada, AR-C467-ahdara-waste

\chapter{Protecting nature}
\label{ch:ar-protecting-nature}
\hlchaptermodality{467}

\begin{chapteropening}
\textbf{By the end of this chapter:} \emph{I can say a nature reserve, endangered, to protect, to die out and to squander.}

Complete the last lesson of chapter 467: I can say a nature reserve, endangered, to protect, to die out and to squander.
\end{chapteropening}

\section[\texorpdfstring{\ar{محمية} (maḥmiyya)}{محمية (maḥmiyya)}]{\ar{محمية} (maḥmiyya) — a nature reserve}

\label{lesson:AR-C467-mahmiyya}



\begin{quote}
Before the new one: say the Arabic for a sand dune, then the Arabic for to choke.
\end{quote}

\subsection*{You'll want to know: \ar{محمية}}
\textbf{\ar{محمية}} — \emph{maḥmiyya} — \textquotedblleft{}a nature reserve\textquotedblright{}.

\begin{quote}
\ar{زرنا} \ar{محمية} \ar{للطيور} \ar{على} \ar{الساحل}.
\end{quote}

\begin{quote}
\emph{zurnā maḥmiyya li-ṭ-ṭuyūr ʿalā as-sāḥil.}
\end{quote}

\begin{quote}
We visited a bird reserve on the coast.
\end{quote}

\begin{cousinweb}
\textbf{\ar{محمية}} (\emph{maḥmiyya}) is built on the native Arabic root \emph{ḥ-m-y}, \textquotedblleft{}to protect, to defend\textquotedblright{} (Lane). It is a passive participle: a protected place.
\end{cousinweb}

\begin{grammarlens}[title={What you've built}]
One more word for getting around.
\end{grammarlens}

\subsection*{Your turn}
\begin{itemize}
\raggedright
  \item \emph{Say it:} \emph{maḥmiyya}
  \item \emph{Say it:} \emph{maḥmiyya}, once more
  \item \emph{Say it:} \emph{zurnā maḥmiyya li-ṭ-ṭuyūr ʿalā as-sāḥil.} — We visited a bird reserve on the coast.
  \item \emph{Recall:} say the Arabic for a sand dune, then the Arabic for to choke, then say \emph{maḥmiyya} again
\end{itemize}

\begin{tcolorbox}[breakable,colback=teal!4,colframe=teal!35!black,title={Before you move on}]
What does \ar{محمية} mean? (\textquotedblleft{}A nature reserve\textquotedblright{}.) Say it once more.
\end{tcolorbox}
\section[\texorpdfstring{\ar{مهدد} (muhaddad)}{مهدد (muhaddad)}]{\ar{مهدد} (muhaddad) — endangered; threatened}

\label{lesson:AR-C467-muhaddad}



\begin{quote}
Before the new one: say the Arabic for to choke, then the Arabic for a nature reserve.
\end{quote}

\subsection*{You'll want to know: \ar{مهدد}}
\textbf{\ar{مهدد}} — \emph{muhaddad} — \textquotedblleft{}endangered; threatened\textquotedblright{}.

\begin{quote}
\ar{النمر} \ar{حيوان} \ar{مهدد} \ar{في} \ar{هذه} \ar{المنطقة}.
\end{quote}

\begin{quote}
\emph{an-namir ḥayawān muhaddad fī hādhihi al-minṭaqa.}
\end{quote}

\begin{quote}
The tiger is an endangered animal in this area.
\end{quote}

\begin{cousinweb}
\textbf{\ar{مهدد}} (\emph{muhaddad}) is built on the native Arabic root \emph{h-d-d}, \textquotedblleft{}to threaten; to pull down\textquotedblright{} (Lane). It is the Form II passive participle, \textquotedblleft{}threatened\textquotedblright{}.
\end{cousinweb}

\begin{grammarlens}[title={What you've built}]
One more word for getting around.
\end{grammarlens}

\subsection*{Your turn}
\begin{itemize}
\raggedright
  \item \emph{Say it:} \emph{muhaddad}
  \item \emph{Say it:} \emph{muhaddad}, once more
  \item \emph{Say it:} \emph{an-namir ḥayawān muhaddad fī hādhihi al-minṭaqa.} — The tiger is an endangered animal in this area.
  \item \emph{Recall:} say the Arabic for to choke, then the Arabic for a nature reserve, then say \emph{muhaddad} again
\end{itemize}

\begin{tcolorbox}[breakable,colback=teal!4,colframe=teal!35!black,title={Before you move on}]
What does \ar{مهدد} mean? (\textquotedblleft{}Endangered; threatened\textquotedblright{}.) Say it once more.
\end{tcolorbox}
\section[\texorpdfstring{\ar{حافظ} (ḥāfaẓa)}{حافظ (ḥāfaẓa)}]{\ar{حافظ} (ḥāfaẓa) — to protect; to preserve}

\label{lesson:AR-C467-hafaza}



\begin{quote}
Before the new one: say the Arabic for a nature reserve, then the Arabic for endangered.
\end{quote}

\subsection*{You'll want to know: \ar{حافظ}}
\textbf{\ar{حافظ}} — \emph{ḥāfaẓa} — \textquotedblleft{}to protect; to preserve\textquotedblright{}.

\begin{quote}
\ar{يحافظ} \ar{الناس} \ar{هنا} \ar{على} \ar{الغابة}.
\end{quote}

\begin{quote}
\emph{yuḥāfiẓu an-nās hunā ʿalā al-ghāba.}
\end{quote}

\begin{quote}
People here look after the forest.
\end{quote}

Present: \ar{يحافظ}, \emph{yuḥāfiẓ}.

It takes \emph{ʿalā}: \emph{ḥāfaẓa ʿalā}, \textquotedblleft{}to take care of\textquotedblright{}.

\begin{cousinweb}
\textbf{\ar{حافظ}} (\emph{ḥāfaẓa}) is built on the native Arabic root \emph{ḥ-f-ẓ}, \textquotedblleft{}to guard, to keep\textquotedblright{} (Lane). It is Form III.
\end{cousinweb}

\begin{grammarlens}[title={What you've built}]
One more word for getting around.
\end{grammarlens}

\subsection*{Your turn}
\begin{itemize}
\raggedright
  \item \emph{Say it:} \emph{ḥāfaẓa}
  \item \emph{Say it:} \emph{ḥāfaẓa}, once more
  \item \emph{Say it:} \emph{yuḥāfiẓu an-nās hunā ʿalā al-ghāba.} — People here look after the forest.
  \item \emph{Recall:} say the Arabic for a nature reserve, then the Arabic for endangered, then say \emph{ḥāfaẓa} again
\end{itemize}

\begin{tcolorbox}[breakable,colback=teal!4,colframe=teal!35!black,title={Before you move on}]
What does \ar{حافظ} mean? (\textquotedblleft{}To protect; to preserve\textquotedblright{}.) Say it once more.
\end{tcolorbox}
\section[\texorpdfstring{\ar{انقرض} (inqaraḍa)}{انقرض (inqaraḍa)}]{\ar{انقرض} (inqaraḍa) — to die out; to become extinct}

\label{lesson:AR-C467-inqarada}



\begin{quote}
Before the new one: say the Arabic for endangered, then the Arabic for to protect.
\end{quote}

\subsection*{You'll want to know: \ar{انقرض}}
\textbf{\ar{انقرض}} — \emph{inqaraḍa} — \textquotedblleft{}to die out; to become extinct\textquotedblright{}.

\begin{quote}
\ar{انقرض} \ar{هذا} \ar{الحيوان} \ar{قبل} \ar{قرون}.
\end{quote}

\begin{quote}
\emph{inqaraḍa hādhā al-ḥayawān qabla qurūn.}
\end{quote}

\begin{quote}
This animal died out centuries ago.
\end{quote}

Present: \ar{ينقرض}, \emph{yanqariḍ}.

\begin{cousinweb}
\textbf{\ar{انقرض}} (\emph{inqaraḍa}) is built on the native Arabic root \emph{q-r-ḍ}, \textquotedblleft{}to cut, to gnaw\textquotedblright{} (Lane). It is Form VII, \textquotedblleft{}to be cut off\textquotedblright{}.
\end{cousinweb}

\begin{grammarlens}[title={What you've built}]
One more word for getting around.
\end{grammarlens}

\subsection*{Your turn}
\begin{itemize}
\raggedright
  \item \emph{Say it:} \emph{inqaraḍa}
  \item \emph{Say it:} \emph{inqaraḍa}, once more
  \item \emph{Say it:} \emph{inqaraḍa hādhā al-ḥayawān qabla qurūn.} — This animal died out centuries ago.
  \item \emph{Recall:} say the Arabic for endangered, then the Arabic for to protect, then say \emph{inqaraḍa} again
\end{itemize}

\begin{tcolorbox}[breakable,colback=teal!4,colframe=teal!35!black,title={Before you move on}]
What does \ar{انقرض} mean? (\textquotedblleft{}To die out; to become extinct\textquotedblright{}.) Say it once more.
\end{tcolorbox}
\section[\texorpdfstring{\ar{أهدر} (ahdara)}{أهدر (ahdara)}]{\ar{أهدر} (ahdara) — to squander; to waste}

\label{lesson:AR-C467-ahdara-waste}



\begin{quote}
Before the new one: say the Arabic for to protect, then the Arabic for to die out.
\end{quote}

\subsection*{You'll want to know: \ar{أهدر}}
\textbf{\ar{أهدر}} — \emph{ahdara} — \textquotedblleft{}to squander; to waste\textquotedblright{}.

\begin{quote}
\ar{أهدر} \ar{المصنع} \ar{ماء} \ar{كثيرا}.
\end{quote}

\begin{quote}
\emph{ahdara al-maṣnaʿ māʾan kathīran.}
\end{quote}

\begin{quote}
The factory wasted a lot of water.
\end{quote}

Present: \ar{يهدر}, \emph{yuhdir}.

\begin{cousinweb}
\textbf{\ar{أهدر}} (\emph{ahdara}) is built on the native Arabic root \emph{h-d-r}, \textquotedblleft{}to be of no account\textquotedblright{} (Lane). It is Form IV, \textquotedblleft{}to treat as worthless\textquotedblright{}.
\end{cousinweb}

\begin{grammarlens}[title={What you've built}]
One more word for getting around.
\end{grammarlens}

\subsection*{Your turn}
\begin{itemize}
\raggedright
  \item \emph{Say it:} \emph{ahdara}
  \item \emph{Say it:} \emph{ahdara}, once more
  \item \emph{Say it:} \emph{ahdara al-maṣnaʿ māʾan kathīran.} — The factory wasted a lot of water.
  \item \emph{Recall:} say the Arabic for to protect, then the Arabic for to die out, then say \emph{ahdara} again
\end{itemize}

\begin{tcolorbox}[breakable,colback=teal!4,colframe=teal!35!black,title={Before you move on}]
What does \ar{أهدر} mean? (\textquotedblleft{}To squander; to waste\textquotedblright{}.) Say it once more.
\end{tcolorbox}
